\subsection{Preparation}\label{sec:initial_setup}

In the preparation stage, we generate the optimized LLVM IR representation of the program, which is used consistently throughout all stages of the pipeline (Figure~\ref{fig:nugget_pipeline}). 
This is because Nugget relies on a stable LLVM IR representation of the program for its unit of work, as introduced in Section~\ref{sec:unit_of_work}.

Optimization is important for the program's performance.
LLVM IR level optimizations, such as loop unrolling, can significantly affect control flow.
To maintain consistency, Nugget applies these optimizations during the generation of the base LLVM IR representation in the preparation stage.
In contrast, most backend optimizations, such as replacing scalar instructions with AVX vector instructions, do not alter the structure of the LLVM IR.
While certain backend passes (e.g., late-stage loop unrolling) may modify LLVM IR layout under specific conditions, such passes are not enabled by default in LLVM.

\shepherd{While build-time modes such as LTO and PGO can alter LLVM IR, they are disabled by default.}
\revision{For our experiments, the LLVM IR structure remained stable across optimization levels, CPU-feature flags, and LLVM versions.
}
Therefore, we can apply the backend optimization in the final binary compilation safely.

Only code included in this base LLVM IR representation of the program will be subject to analysis.
For instance, any linked libraries must also be processed through Nugget during this preparation stage; otherwise, their contents will be excluded from the analysis.
This selective inclusion also allows us to ignore irrelevant or non-contributory operations, such as spin loops in multi-threaded libraries~\cite{looppoint}, that do not meaningfully impact program progress.
